\documentclass[11pt]{article}
\usepackage[T1]{fontenc}
\usepackage[utf8]{inputenc}
\usepackage{lmodern}
\usepackage[margin=1in]{geometry}
\usepackage{amsmath,amssymb,amsthm,mathtools}
\usepackage{booktabs,array}
\usepackage{enumitem}
\usepackage[numbers,sort&compress]{natbib}
\usepackage{microtype}
\usepackage[colorlinks=true,allcolors=blue]{hyperref}
\usepackage[nameinlink,noabbrev]{cleveref}
\newtheorem{theorem}{Theorem}[section]
\newtheorem{proposition}[theorem]{Proposition}
\newtheorem{lemma}[theorem]{Lemma}
\newtheorem{corollary}[theorem]{Corollary}
\theoremstyle{definition}
\newtheorem{definition}[theorem]{Definition}
\newtheorem{example}[theorem]{Example}
\theoremstyle{remark}
\newtheorem{remark}[theorem]{Remark}
\DeclareMathOperator{\conv}{conv}
\DeclareMathOperator{\rank}{rank}
\DeclareMathOperator{\supp}{supp}
\DeclareMathOperator{\aff}{aff}
\DeclareMathOperator{\lin}{lin}
\newcommand{\R}{\mathbb R}
\newcommand{\Q}{\mathbb Q}
\newcommand{\Z}{\mathbb Z}
\newcommand{\one}{\mathbf 1}
\newcommand{\Obs}{\mathcal O}
\newcommand{\Blocks}{\mathcal B}
\newcommand{\Hull}{H}
\newcommand{\Flow}{P}
\newcommand{\states}{\{0,\ldots,m\}}
\title{Sparse convex hulls for network flows coupled to a simplex}
\author{}
\date{}
\begin{document}
\maketitle
\begin{abstract}
We study the convex hull of selected bilinear products between a bounded
equality-constrained network flow and a simplex vector. The network and the
retained product coordinates provide two distinct sources of structure:
circulations factor across undirected blocks, and each block need retain only
the simplex states observed there. We develop exact formulations and
original-coordinate descriptions from these reductions, and examine the
limits of small-coefficient descriptions. The starting point is classical
simplex-vertex disaggregation; the structural question is how much of that
formulation is needed for a sparse set of products.
\end{abstract}

\section{Introduction}\label{sec:introduction}
Network models often combine continuous routing decisions with a choice among
mutually exclusive services or operating modes. A route-specific adjustment
to the cost of a service gives a product of an arc flow and a service weight.
Only a small fraction of the possible route--service pairs may occur in the
model. Convexifying each product independently discards restrictions linking
flows through the network and linking service weights through the simplex.
An exact joint hull retains these restrictions, but its standard extended
formulation introduces a copy of every arc flow for every simplex vertex.

This paper studies the structure of that joint hull when only selected
products are retained. We distinguish the size of an extended formulation
from the coefficients and separation cost of an original-coordinate
description. We also distinguish equality-constrained flow models from
network models with balance inequalities: adding slack arcs to convert the
latter to the former can change the graph structure on which a specialized
result depends. The elementary foundations below make the state and graph
decompositions explicit, including degenerate capacities and zero simplex
weights.

\section{Model, prior work, and disaggregation}\label{sec:foundations}

\subsection{The sparse network--simplex hull}\label{subsec:model}
Let $D=(V,E)$ be a finite directed multigraph. Parallel arcs and self-loops
are permitted. Its incidence matrix $A\in\{-1,0,1\}^{V\times E}$ has a $-1$
at the tail and a $+1$ at the head of a non-loop arc; a self-loop has a zero
column. Thus $Ax=b$ uses the convention \emph{incoming minus outgoing}.
For $b\in\Q^V$ and finite $u\in\Q_+^E$, set
\begin{equation}\label{eq:flow-polytope}
 \Flow=\{x\in\R^E:Ax=b,\ 0\le x\le u\}.
\end{equation}
Lower bounds can be handled by translating $x$ and the corresponding product
coordinates. Throughout, graph assumptions concern the graph in
\eqref{eq:flow-polytope} after any such modeling transformations.

For an integer $m\ge0$, write $[m]=\{1,\ldots,m\}$ and
\begin{equation}\label{eq:simplex}
 \Delta_m=\{y\in\R_+^m:\one^\top y\le1\},\qquad
 \lambda_j=y_j\ (j\in[m]),\qquad
 \lambda_0=1-\one^\top y.
\end{equation}
The \emph{states} are the $d=m+1$ vertices $e^0=0,e^1,\ldots,e^m$ of
$\Delta_m$; $e^j$ is the $j$th unit vector when $j\ge1$. The residual state
is state $0$. An observation set $\Obs\subseteq E\times[m]$ specifies
the products retained as original coordinates. We study
\begin{align}
 S&=\{(x,y,z):x\in\Flow,\ y\in\Delta_m,
             \ z_{ej}=x_e y_j\quad((e,j)\in\Obs)\},\label{eq:graph-set}\\
 \Hull&=\conv S.\label{eq:sparse-hull}
\end{align}
An unobserved product is absent from $z$, even if it is introduced temporarily
in an extended formulation. State $0$ has no observed coordinates.
We use $\conv\varnothing=\varnothing$. When $m=0$, $y$ and $z$ are empty
vectors and $\Hull=\Flow$.

The original coordinates are $(x,y,z)$. An extended formulation is a linear
system in these coordinates and additional continuous variables whose
projection is $\Hull$. A separation procedure must either certify membership
or produce a linear inequality valid for all of $\Hull$ and strictly violated
by the candidate. Unless specified otherwise, a facet is taken relative to
the affine hull. A bound on integer coefficients presupposes a stated
normalization; a ratio of two nonzero coefficients is invariant under scaling
the inequality. Arithmetic-operation bounds are distinguished from bounds
on rational encoding length.

For comparison, the usual independent product relaxation adds, for every
$(e,j)\in\Obs$, the four inequalities
\begin{equation}\label{eq:mccormick}
 0\le z_{ej},\qquad z_{ej}\le u_e y_j,\qquad
 z_{ej}\le x_e,\qquad
 z_{ej}\ge x_e-u_e(1-y_j)
\end{equation}
to $x\in\Flow$ and $y\in\Delta_m$. These are the McCormick inequalities
for the box $[0,u_e]\times[0,1]$ \citep{McCormick1976}. They need not describe
the joint hull in \eqref{eq:sparse-hull}.

\subsection{Relation to prior work}\label{subsec:prior-work}
The general convexification principle is established. Disjunctive hull
formulations \citep{Balas1998} apply after expanding the simplex into its
vertices. Simultaneous convexification over a polytope times a simplex is
developed by \citet{DavarniaRichardTawarmalani2017}.
\citet{KhademniaDavarnia2025} specialize this setting to network polytopes:
their electronic companion, p.~1, equation~(25), gives a full extended
hull,\footnote{The published electronic companion is available at
\href{https://pubsonline.informs.org/doi/suppl/10.1287/moor.2023.0001/suppl_file/moor.2023.0001.sm1.pdf}{this direct supplement link}.}
and their
aggregation-and-projection framework gives a complete original-space
description in principle. Their explicit tree and forest constructions
identify more structured inequality families; the forest recipe for multiple
simplex coordinates should not be confused with the entire projection
framework. Their general framework permits balance inequalities, and their
Section~3 represents equality balances by pairs of opposite inequalities.
Thus their network specialization includes \eqref{eq:flow-polytope}; here we
exploit its underlying equality circulation space. We use the published
article and its electronic companion for equation and example numbering.

The shared-simplex gluing principle is also known. In particular,
\citet[Proposition~2.6]{Davarnia2016} proves that convexifying Cartesian
components separately is exact when the components share the same simplex.
The cited locator belongs to the dissertation, which is distinct from the
2017 journal article. Our block decomposition applies this principle to the
independent circulation spaces of an equality-constrained network. We give a
direct proof to make the treatment of sparse observations and zero weights
explicit.

Several established methods explain the subsequent structural reductions.
Reduced reformulation--linearization uses multiplied linear equations to
remove common-factor products \citep{LibertiPantelides2006}. In networks,
the remaining freedom can be identified with circulations supported on
unobserved arcs. The network representation of Cayley hulls of
\citet{KisHorvath2022} gives cut descriptions and transportation projections;
their Section~5.7 is a close predecessor for the parallel-path construction.
The support geometry of Minkowski sums and edge-generated zonotope
refinements is classical \citep{GritzmannSturmfels1993,OnnRothblum2004}.
These are ingredients of the graph-specific formulations and separation
procedures, rather than new general convexification mechanisms.

Two further comparisons delimit coefficient obstructions. Three-way
transportation universality is established by \citet{DeLoeraOnn2006}, and
large coefficients in general $0/1$ polytopes are classical
\citep{AlonVu1997}. A restricted network realization must therefore be
specified as part of a coefficient result. Moreover, the integrality of
\emph{aggregate} flows on two-terminal series--parallel networks does not
give a hull description in selected state-flow coordinates.
\citet[Remark~1]{AlmoghrabiSkutellaWarode2026} explicitly distinguish
aggregate integrality from decomposability of the vector of individual
commodity flows. Our observed products retain such individual-state
information. Neither a sparse extended formulation nor polynomial separation
precludes large numerical coefficients after projection.

\subsection{The known exact disaggregation}\label{subsec:disaggregation}
For $\lambda\ge0$, define the scaled flow system
\begin{equation}\label{eq:scaled-flow}
 \Flow(\lambda)=\{f\in\R^E:Af=\lambda b,\ 0\le f\le\lambda u\}.
\end{equation}
For $\lambda>0$, this is $\lambda\Flow$; at $\lambda=0$ it is $\{0\}$,
including when $\Flow$ is empty. Using \eqref{eq:scaled-flow} avoids assigning
an ambiguous meaning to $0\varnothing$.

\begin{proposition}[Simplex-vertex disaggregation]\label{prop:disaggregation}
A point $(x,y,z)$ belongs to $\Hull$ if and only if $y\in\Delta_m$ and
there exist state flows $f^0,\ldots,f^m\in\R^E$ such that
\begin{equation}\label{eq:disaggregation}
 f^k\in\Flow(\lambda_k)\quad(k\in\states),\qquad
 \sum_{k=0}^m f^k=x,\qquad
 f^j_e=z_{ej}\quad((e,j)\in\Obs).
\end{equation}
The equivalence holds even if $\Flow$ is empty. If the system is feasible,
it provides a convex decomposition into at most $m+1$ points of $S$.
\end{proposition}
\begin{proof}
For a point of $S$, set $f^k=\lambda_k x$. This satisfies all displayed
relations, since $\sum_k\lambda_k=1$. The system is linear in
$(x,y,z,f)$, so its projection is convex and contains $\Hull$.

Conversely, suppose the system is feasible. If $\lambda_k=0$, the finite
upper bounds give $f^k=0$. For each positive weight let
$x^k=f^k/\lambda_k\in\Flow$ and define the graph point
$p^k=(x^k,e^k,z^k)$ by
$z^k_{ej}=x^k_e(e^k)_j$. Its weighted average has $x$ coordinate
$\sum_k f^k=x$, $y$ coordinate $\sum_k\lambda_k e^k=y$, and observed
product coordinate $\lambda_j x^j_e=f^j_e=z_{ej}$ when
$\lambda_j>0$. If $\lambda_j=0$, that observed coordinate is zero on both
sides. Thus $(x,y,z)=\sum_{k:\lambda_k>0}\lambda_k p^k\in\Hull$.
There is at least one positive weight, so feasibility also implies
$\Flow\ne\varnothing$. This proves the empty-set assertion.
\end{proof}

The formulation has $(m+1)|E|$ state-flow variables before elimination.
Its rational size is polynomial in the ordinary network and simplex input,
so generic polynomial-time hull optimization and separation already follow
from linear programming. The structural issue is avoiding this full
state-by-arc expansion, or projecting it into useful original-coordinate
inequalities. In particular, eliminating $f^0=x-\sum_{j=1}^m f^j$ keeps
the residual bounds
\begin{equation}\label{eq:residual-bounds}
 0\le x-\sum_{j=1}^m f^j\le
       \Bigl(1-\sum_{j=1}^m y_j\Bigr)u.
\end{equation}
These bounds cannot in general be discarded.

\section{Independent blocks and locally merged states}\label{sec:blocks}

\subsection{Circulation factorization}\label{subsec:block-factorization}
All graph connectivity terminology below refers to the undirected multigraph
underlying $D$. Isolated vertices impose only their balance equation. Let
$E_{\rm br}$ be the bridges. Let $\Blocks$ be the cyclic blocks: the maximal
biconnected edge blocks that contain a cycle, with each self-loop treated as
its own block. A pair of parallel edges can form a cycle of length two.
The sets $E_{\rm br}$ and $E_B$ for $B\in\Blocks$ partition $E$.
The cycle rank of a connected block is
\begin{equation}\label{eq:cycle-rank}
 r_B=|E_B|-|V_B|+1.
\end{equation}
In particular, a loop has rank one. Rank refers to the ambient circulation
space, even when capacities reduce the dimension of the feasible polytope.

The equation $Av=b$ is solvable exactly when the entries of $b$ sum to zero
on every connected component. Necessity follows by summing the incidence
rows in a component. For sufficiency, choose a spanning forest, set the
nonforest coordinates to zero, and determine each tree-edge flow by the
balance sum on one side of its fundamental cut, with the sign determined by
the arc orientation. This satisfies every balance equation. This construction
uses $O(|V|+|E|)$ arithmetic operations, and for rational $b$ produces a
rational $v$ of polynomial encoding length. The reference flow $v$ need not
obey its capacity bounds. If the component condition fails, $\Flow$ and
$\Hull$ are empty.

Fix such a reference vector $v$. For each $B$, let $A_B$ be the incidence
matrix restricted to columns in $E_B$ and rows in $V_B$, and define
\begin{equation}\label{eq:block-domain}
 K_B=\{h\in\R^{E_B}:A_Bh=0,\ -v_{E_B}\le h\le u_{E_B}-v_{E_B}\}.
\end{equation}
Unlike $\Flow$, this is a domain of circulation \emph{deviations}; it need
not contain the origin.

\begin{lemma}[Block factorization]\label{lem:block-factorization}
Every $g\in\ker A$ vanishes on $E_{\rm br}$ and restricts to a circulation
on each cyclic block. Conversely, combining arbitrary block circulations
and zero bridge coordinates gives an element of $\ker A$. Consequently,
if every bridge satisfies $0\le v_e\le u_e$, the map
\begin{equation}\label{eq:block-product-map}
 x_e=v_e\quad(e\in E_{\rm br}),\qquad
 x_{E_B}=v_{E_B}+h_B\quad(B\in\Blocks)
\end{equation}
is an affine bijection from $\prod_{B\in\Blocks}K_B$ to $\Flow$.
If a bridge bound fails, $\Flow$ is empty. The product of an empty family
of block domains is a singleton.
\end{lemma}
\begin{proof}
Choose a spanning forest of the underlying multigraph, omitting loops. Each
nonforest edge gives a signed fundamental-cycle circulation: use the unique
forest path between its ends, or just the edge itself for a loop. These
vectors form a basis of $\ker A$. Indeed, subtracting from a circulation
the combination that cancels all nonforest coordinates leaves a circulation
supported on a forest, which is zero by successive leaf elimination.

Each fundamental cycle lies in a single cyclic block. Hence its vector is
zero on every bridge, and on every block it is either zero or a circulation
of that block. The basis expansion proves the first assertion. For the
converse, extend each block circulation by zero to $E$. Its incidence is
zero at every vertex, including articulation vertices, so their sum is a
global circulation. The sum is direct because the block edge sets are
disjoint. Applying this decomposition to $x-v$ proves that $Ax=b$ is
equivalent to the circulation conditions and the fixed bridge coordinates
in \eqref{eq:block-product-map}. The bounds then separate by edge block,
giving precisely \eqref{eq:block-domain} and the bridge conditions.
\end{proof}

In particular, an articulation vertex does not carry an additional coupling
variable between block deviations. Its fixed net balance is already carried
by $v$; every block deviation has zero net balance at that vertex separately.
Any spanning tree in $B$ gives a fundamental-cycle matrix
$C_B\in\{0,\pm1\}^{E_B\times r_B}$ with chord rows equal to the identity.
Thus $h_B=C_Bt_B$ gives unique cycle coordinates $t_B\in\R^{r_B}$.
These coordinates remain useful when $K_B$ is lower-dimensional.

\subsection{The state merger and its exact refinement}\label{subsec:state-merger}
For a nonempty convex set $K$, use the usual scalar multiplication
$\alpha K=\{\alpha h:h\in K\}$, so $0K=\{0\}$. For nonnegative scalars,
\begin{equation}\label{eq:homothetic-merger}
 \alpha_1K+\cdots+\alpha_qK
       =(\alpha_1+\cdots+\alpha_q)K.
\end{equation}
If the total weight is positive, the inclusion from left to right follows
by dividing by that weight and taking a convex combination; the reverse
inclusion uses the same point of $K$ in each summand. If the total is zero,
both sides are $\{0\}$. An empty sum is also interpreted as $\{0\}$.

For each block define its observed labels and their count by
\begin{equation}\label{eq:observed-labels}
 J_B=\{j\in[m]: (e,j)\in\Obs\text{ for some }e\in E_B\},
 \qquad a_B=|J_B|.
\end{equation}
This is a structural definition: a label in $J_B$ remains present when
$y_j=0$. Let $*_B$ denote a local merged state, distinct from the global
residual state $0$, and set
\begin{equation}\label{eq:merged-weight}
 I_B=J_B\cup\{*_B\},\qquad
 \lambda_{B,j}=y_j\ (j\in J_B),\qquad
 \lambda_{B,*}=1-\sum_{j\in J_B}y_j.
\end{equation}
The merged state collects all labels absent from $J_B$ and state $0$.
Different blocks may merge different global states.

\begin{theorem}[Exact blockwise state reduction]\label{thm:block-state-reduction}
Assume $\Flow\ne\varnothing$. A point $(x,y,z)$ belongs to $\Hull$ if and
only if $x\in\Flow$, $y\in\Delta_m$, every observed bridge satisfies
\begin{equation}\label{eq:bridge-product}
 z_{ej}=v_e y_j\qquad(e\in E_{\rm br},\ (e,j)\in\Obs),
\end{equation}
and for every cyclic block $B$ there exist vectors
$h^{B,k}\in\R^{E_B}$, $k\in I_B$, satisfying
\begin{align}
 h^{B,k}&\in\lambda_{B,k}K_B\quad(k\in I_B),\label{eq:local-state-domain}\\
 \sum_{k\in I_B}h^{B,k}&=x_{E_B}-v_{E_B},\label{eq:local-state-sum}\\
 h^{B,j}_e&=z_{ej}-y_jv_e
       \quad((e,j)\in\Obs,\ e\in E_B).\label{eq:local-observation}
\end{align}
The block conditions yield state flows in
\eqref{eq:disaggregation}, and hence a decomposition into at most $m+1$
graph points, by explicit proportional refinement.
\end{theorem}
\begin{proof}
By \cref{lem:block-factorization}, every $K_B$ is nonempty. Start with
state flows in \eqref{eq:disaggregation}. Their deviations
$g^k=f^k-\lambda_k v$ are circulations. On $B$, they satisfy
$g^k_{E_B}\in\lambda_kK_B$: division proves this for positive weights,
and both deviations and scaled domains are zero for zero weights.
For $j\in J_B$ take $h^{B,j}=g^j_{E_B}$, and sum the other deviations
into $h^{B,*}$. Identity \eqref{eq:homothetic-merger} proves
\eqref{eq:local-state-domain}. Summing the deviations proves
\eqref{eq:local-state-sum}, and the observed coordinates give
\eqref{eq:local-observation}. On bridges the deviations vanish, proving
\eqref{eq:bridge-product}.

Conversely, fix feasible local vectors. For each global state $k$ and block
$B$, if $k\in J_B$, set $g^k_{E_B}=h^{B,k}$. For the remaining states use
\begin{equation}\label{eq:proportional-refinement}
 g^k_{E_B}=\begin{cases}
 (\lambda_k/\lambda_{B,*})h^{B,*},&\lambda_{B,*}>0,\\
 0,&\lambda_{B,*}=0.
 \end{cases}
\end{equation}
If the merged weight is zero, every constituent weight is zero as well.
In every case $g^k_{E_B}\in\lambda_kK_B$, and the refined vectors sum to
the block deviation. Put $g^k_e=0$ on bridges and combine blocks. The block
factorization gives $Ag^k=0$. Then $f^k=\lambda_kv+g^k$ satisfies
$Af^k=\lambda_kb$. Its capacity bounds hold on blocks by the scaled
domains, and on bridges because $\Flow\ne\varnothing$ implies
$0\le v_e\le u_e$. Equations \eqref{eq:local-state-sum} and
\eqref{eq:local-observation}, together with the bridge equations, give
$\sum_kf^k=x$ and all observed products. Apply
\cref{prop:disaggregation}.
\end{proof}

The domain condition \eqref{eq:local-state-domain} is itself linear:
\begin{equation}\label{eq:linear-scaled-block}
 A_Bh^{B,k}=0,\qquad
 -\lambda_{B,k}v_{E_B}\le h^{B,k}
       \le\lambda_{B,k}(u_{E_B}-v_{E_B}).
\end{equation}
Finite bounds force $h^{B,k}=0$ at zero weight. No division by a variable is
needed in the formulation; division appears only in the recovery proof at
strictly positive weights. If $J_B=\varnothing$, the block has one merged
state of weight one and its conditions reduce to $x_{E_B}-v_{E_B}\in K_B$,
already enforced by $x\in\Flow$. Thus unobserved blocks require no new
variables or constraints.

The theorem is a constructive specialization of the known shared-simplex
gluing principle. It does not introduce independent simplex vectors for
different blocks: all use the same global $y$. It also does not merge
observed state slices, whose domains can differ after the products are fixed.

\begin{example}[A common matrix is not sufficient for merging]\label{ex:merger-boundary}
Let $M=(-1,1,2)^\top$, $c^1=(0,1,100)^\top$, and
$c^2=(0,100,2)^\top$. Both polytopes
$Q_i=\{q\in\R:Mq\le c^i\}$ equal $[0,1]$.
For weights $\alpha_1=\alpha_2=1/2$, the weighted Minkowski sum is still
$[0,1]$, whereas the aggregated rows
$Mq\le(c^1+c^2)/2$ allow $q=2$. The two upper bounds become
$q\le101/2$ and $2q\le51$. Thus replacing state-specific systems by
the sum of their right-hand sides is generally inexact, even for a common
matrix and nonempty bounded domains. This is the standard reaggregation
boundary discussed by \citet[Section~2, equation~(7)]{KisHorvath2022}.
The sets in this example are identical, so the set identity
\eqref{eq:homothetic-merger} still holds. What fails is its replacement by
addition of the different, redundant right-hand-side descriptions. In
\cref{thm:block-state-reduction}, every state instead uses the single fixed
representation \eqref{eq:block-domain}, with its bounds scaled by the state
weight as in \eqref{eq:linear-scaled-block}. Adding these scaled bounds gives
that same representation scaled by the total weight, so the set identity
justifies the merged formulation.
\end{example}

Finally, these results concern the hull of \eqref{eq:graph-set}. Intersecting
$\Hull$ with additional linear side constraints yields a valid relaxation
of the graph with those constraints, but need not give its exact convex
hull. Such constraints can couple blocks or alter conditional state domains;
exactness then requires a separate argument.

\section{Compression by cycles and observations}\label{sec:compression}
This section compresses the exact block formulation in two steps. First,
cycle coordinates replace arc coordinates and suppress long paths. Second,
observations determine some cycle coordinates, leaving only cycles supported
on unobserved arcs. We assume $\Flow\ne\varnothing$; if the original flow
domain is empty, its hull is described by any inconsistent linear system.

\subsection{Fixed coordinates and reduced graphs}\label{subsec:fixed-arcs}
The structural parameters can be applied after elementary exact preprocessing.
This is useful when capacities force a nominal cycle to carry no freedom.

\begin{lemma}[Deleting fixed arcs]\label{lem:fixed-arcs}
Let $F\subseteq E$ and suppose $x_e=c_e$ throughout $\Flow$ for every
$e\in F$, where $c\in\Q^F$ is known. Set $R=E\setminus F$,
$b'=b-A_{:,F}c$, and
\[
 \Flow'=\{x_R:A_{:,R}x_R=b',\ 0\le x_R\le u_R\},\qquad
 \Obs'=\Obs\cap(R\times[m]).
\]
The hull for $(\Flow',\Obs')$, together with
$x_F=c$ and $z_{ej}=c_e y_j$ for observed $e\in F$, is affinely
isomorphic to $\Hull$. If $F$ contains every coordinate fixed throughout
$\Flow$, then
\begin{equation}\label{eq:reduced-affine-hull}
 \aff\Flow'=\{x_R:A_{:,R}x_R=b'\}.
\end{equation}
\end{lemma}
\begin{proof}
Feasible vectors of $\Flow$ restrict to $\Flow'$. Conversely, adjoining
$c$ to a vector of $\Flow'$ satisfies the original balances and bounds,
because $0\le c\le u_F$. The same bijection at graph points adjoins the
affine coordinates $z_{ej}=c_e y_j$, so it commutes with convexification.

For the last assertion, every remaining coordinate is nonconstant on
$\Flow'$. For each $e\in R$, choose two feasible vectors differing in
coordinate $e$; their midpoint has $0<x_e<u_e$. Average these midpoints
over $e\in R$. The resulting feasible vector is strictly inside every
remaining bound. A sufficiently small relative neighborhood of that vector
in the balance affine space therefore remains feasible, proving
\eqref{eq:reduced-affine-hull}. If $R$ is empty, both spaces are the same
singleton and the assertion is immediate.
\end{proof}

The lemma takes the fixed coordinates as input; it does not require a
particular detection algorithm. Deleting them can split blocks or reduce
their cycle ranks. Once all fixed coordinates are removed, the remaining
cycle space is the direction space of the feasible affine hull. Conditional
slices at a particular observation or zero state weight can still have
smaller dimension.

\subsection{Suppressing paths}\label{subsec:suppression}
Fix a cyclic block $B$. Degrees in this subsection are degrees \emph{within
the block}, with loops counted twice. For $r_B\ge2$, suppress maximal paths
whose internal vertices have degree two, and give each resulting core edge
an arbitrary orientation. A rank-one block is a single cycle; represent
its entire cyclic traversal by one core loop. This includes two parallel
edges and an original self-loop.

Write $\mathcal P_B$ for the core edges, or equivalently the original
paths, and $k_B=|\mathcal P_B|$. For an original arc $e$ on path $p$, let
$\varepsilon_e\in\{-1,1\}$ compare its orientation with the chosen
path traversal. A circulation has constant signed deviation on a path:
\begin{equation}\label{eq:path-deviation}
 x_e=v_e+\varepsilon_e\delta_p\qquad(e\in p).
\end{equation}
The reason is balance of the block circulation at each internal degree-two
vertex. Balances at the unsuppressed vertices say precisely that $\delta$
is a circulation of the core. Arc bounds give the path interval
\begin{align}
 L_p&=\max_{e\in p}\ell_e,& U_p&=\min_{e\in p}u'_e,\label{eq:path-intervals}\\
 [\ell_e,u'_e]&=
 \begin{cases}[-v_e,u_e-v_e],&\varepsilon_e=1,\\
 [v_e-u_e,v_e],&\varepsilon_e=-1.
 \end{cases}\nonumber
\end{align}
All intervals are finite. Their joint feasibility follows from
$\Flow\ne\varnothing$.

Suppression preserves cycle rank. If $r_B\ge2$, the connected core has
minimum degree at least three: otherwise the original block would have
been a single cycle. If its vertex count is $n_B$, then
$2k_B\ge3n_B$ and $r_B=k_B-n_B+1$, so
\begin{equation}\label{eq:core-size}
 n_B\le2r_B-2,\qquad k_B\le3r_B-3.
\end{equation}
For rank one, $k_B=1$. In particular, $k_B\le3r_B$ in every case.

Choose a spanning tree of the core and let
$\widehat C_B\in\{0,\pm1\}^{k_B\times r_B}$ be its fundamental-cycle
matrix, normalized to have identity chord rows. Its columns parametrize
core circulations bijectively. Thus the feasible core domain is
\begin{equation}\label{eq:core-domain}
 T_B=\{t\in\R^{r_B}:L_p\le(\widehat C_Bt)_p\le U_p
                         \quad(p\in\mathcal P_B)\}.
\end{equation}
It is nonempty and bounded, including in lower-dimensional cases. Boundedness
follows already from the chord rows. For later use, choose one original
representative arc $e(p)$ on each path and put
\begin{equation}\label{eq:aggregate-path-expression}
 s_p(x)=\varepsilon_{e(p)}(x_{e(p)}-v_{e(p)}),\qquad
 q_{ej}=\varepsilon_e(z_{ej}-v_e y_j).
\end{equation}
For $x\in\Flow$, $s_p(x)$ is its core deviation. A separate observation
$q_{ej}$ is retained for \emph{every} observed original arc, including
multiple observations of one path.

\begin{theorem}[Cycle-coordinate extended hull]\label{thm:cycle-compression}
The hull $\Hull$ has an exact rational extended formulation with
\begin{equation}\label{eq:initial-variable-count}
 \sum_{B\in\Blocks}r_Ba_B
\end{equation}
additional variables and
\begin{equation}\label{eq:compression-row-count}
 O\left(|V|+|E|+m+|\Obs|
       +\sum_{B:a_B>0}r_B(a_B+1)\right)
\end{equation}
rows, counting the original domain constraints. Coefficients of $x$, $z$,
and additional variables can all be chosen from $\{0,\pm1\}$. Coefficients
of $y$ and the right-hand sides may depend on the rational network data.
All coefficient encoding lengths are polynomial in the input length.
\end{theorem}
\begin{proof}
Retain $x\in\Flow$, $y\in\Delta_m$, and the bridge equations
\eqref{eq:bridge-product}. For each $B$ with $a_B>0$, introduce
$t^{B,j}\in\R^{r_B}$ for $j\in J_B$ and impose
\begin{align}
 y_jL_p&\le(\widehat C_Bt^{B,j})_p\le y_jU_p,
       &&j\in J_B,\ p\in\mathcal P_B,\label{eq:compressed-state}\\
 q_{ej}&=(\widehat C_Bt^{B,j})_p,
       &&(e,j)\in\Obs,\ e\in p\subseteq E_B,\label{eq:compressed-observation}\\
 \lambda_{B,*}L_p&\le s_p(x)-\sum_{j\in J_B}
                      (\widehat C_Bt^{B,j})_p
                    \le\lambda_{B,*}U_p,
       &&p\in\mathcal P_B.\label{eq:compressed-residual}
\end{align}
The aggregate deviation $s(x)$ is a core circulation. Hence the middle
quantity in \eqref{eq:compressed-residual} is a core circulation as well.
These constraints are exactly \cref{thm:block-state-reduction} in core
coordinates, with the merged state eliminated by its sum equation.
At a zero weight, the finite path bounds force every core path value to
zero; the identity chord rows force the entire state vector to zero.
Consequently the equivalence includes all simplex boundary strata. Blocks
with $a_B=0$ add no constraints, as in the block theorem.

There are $r_Ba_B$ variables, $2k_B(a_B+1)$ bound rows, and one equation
per block observation. Bound \eqref{eq:core-size} proves the row count.
Every state variable has coefficient $0,1$, or $-1$. An observation row
has one original product coordinate, and each residual row uses one
original flow coordinate through $s_p(x)$; no sums of different cycle
expressions for that flow coordinate are necessary. This proves the stated
coefficient bound. Reference flows are sums of rational balances; path
endpoints are selected signed capacity offsets. The remaining coefficients
are sums of these data, with polynomial rational encoding length.
\end{proof}

\subsection{Eliminating observed directions}\label{subsec:observed-directions}
Define the observed and unobserved arcs for a block and label by
\[
 O_{Bj}=\{e\in E_B:(e,j)\in\Obs\},\qquad U_{Bj}=E_B\setminus O_{Bj}.
\]
Let $c(V_B,U_{Bj})$ count all connected components of this unobserved
undirected subgraph, including isolated vertices. Its cycle rank is
\begin{equation}\label{eq:unobserved-rank}
 \rho_{Bj}=|U_{Bj}|-|V_B|+c(V_B,U_{Bj}).
\end{equation}
Loops and parallel edges keep their usual multigraph meanings. In particular,
each unobserved loop contributes one to the cycle rank.

\begin{lemma}[Unit minors for cycle elimination]\label{lem:cycle-tu}
A fundamental-cycle matrix with identity chord rows is totally unimodular.
Let $C$ be such a matrix, let $D$ consist of $d$ linearly independent rows,
and choose $d$ columns $I$ such that $B=D_{:,I}$ is nonsingular. Write $F$
for the complementary columns. Then, for every row $c$ of $C$, the vectors
\begin{equation}\label{eq:unit-schur}
 W=c_I B^{-1},\qquad R=c_F-WD_{:,F}
\end{equation}
have entries in $\{0,\pm1\}$.
\end{lemma}
\begin{proof}
Delete one incidence row of the connected core and partition the reduced
incidence matrix as $[N_T\ N_Q]$ into tree and chord columns. A tree
basis has determinant $\pm1$. Every square incidence minor is $0$ or
$\pm1$: induction expands a column with at most one nonzero entry; if
every column has two entries, they are opposite and the row sum is zero.
Every minor of $N_T^{-1}N_Q$ equals, up to sign, the determinant of a
matrix obtained by replacing the corresponding tree columns by chord
columns, divided by $\det N_T$. This follows by multiplying the replacement
matrix by $N_T^{-1}$ and expanding its unchanged identity columns.
The replacement is a square incidence submatrix. Therefore
$-N_T^{-1}N_Q$ is totally unimodular. Appending identity chord rows
preserves this property, proving the assertion for $C$. A core consisting
of one loop has $C=(1)$ directly.

In particular $\det B=\pm1$. Cramer's rule applied to $WB=c_I$ expresses
each entry of $W$ as a selected-row replacement minor of $C$ divided by
$\det B$. For a free column $f$, the bordered determinant identity gives
\[
 R_f=\frac{\det\begin{pmatrix}B&D_{:,f}\\c_I&c_f\end{pmatrix}}
                  {\det B}.
\]
The numerator is a minor of $C$, or zero when its rows repeat. These
formulas prove \eqref{eq:unit-schur}. Empty pivot or free sets have their
usual zero-dimensional matrix interpretation.
\end{proof}

\begin{theorem}[Observation-sensitive extended hull]\label{thm:observed-compression}
There is an exact rational extended formulation for $\Hull$ with
\begin{equation}\label{eq:observed-variable-count}
 \sum_{B\in\Blocks}\ \sum_{j\in J_B}\rho_{Bj}
\end{equation}
additional variables and the row bound \eqref{eq:compression-row-count}.
Coefficients of $x$, $z$, and additional variables belong to $\{0,\pm1\}$;
simplex coefficients and right-hand sides have polynomial rational encoding
length. The variable count is an upper bound for this construction, not a
lower bound for arbitrary extended formulations.
\end{theorem}
\begin{proof}
Fix $B$ and $j\in J_B$, suppress their subscripts, and write
$C=\widehat C_B$, $r=r_B$. Choose a maximal independent set of observed
path rows of $C$, giving $D\in\R^{d\times r}$. For each selected row
choose one actual observed arc and use its value from
\eqref{eq:aggregate-path-expression}, giving $q\in\R^d$. The selected
entries of $q$ involve distinct original product variables. All other
observations will remain as equations, including any repeated path values.

Choose pivot columns $I$ and complementary columns $F$ as in
\cref{lem:cycle-tu}. Retain $g=t_F\in\R^{r-d}$ and eliminate the pivot
coordinates by the reversible substitution
\begin{equation}\label{eq:pivot-reconstruction}
 t_I=B^{-1}(q-D_{:,F}g).
\end{equation}
For every core path $p$ with row $c_p$ define
\begin{equation}\label{eq:path-reconstruction}
 \phi_{pj}=(c_{p,I}B^{-1})q+
            (c_{p,F}-c_{p,I}B^{-1}D_{:,F})g.
\end{equation}
Replace $(\widehat C_Bt^{B,j})_p$ by $\phi_{pj}$ in
\eqref{eq:compressed-state}--\eqref{eq:compressed-residual}. Selected
observation equations become identities and can be omitted; all other
observation equations remain. Formula \eqref{eq:pivot-reconstruction}
proves equivalence in both directions, so no projection condition is lost.

By \cref{lem:cycle-tu}, each selected product and retained variable has
coefficient $0,1$, or $-1$ in every $\phi_{pj}$. A nonselected observation
adds a distinct product coordinate; it cannot add a second occurrence of
a selected product. Different labels use disjoint product and auxiliary
columns. The residual uses one original flow coordinate through $s_p(x)$.
Thus collecting terms does not increase any $x$, $z$, or auxiliary
coefficient beyond one in magnitude. Only coefficients of $y$ collect
reference-flow and interval data. Their encoding lengths remain polynomial.
Elimination adds no rows.

It remains to identify $r-d$. Restrict a block circulation to $O_{Bj}$.
The kernel is precisely the space of circulations supported on $U_{Bj}$,
whose dimension is \eqref{eq:unobserved-rank} by the incidence-rank formula.
The core matrix parametrizes all block circulations, and observing any
arc of a suppressed path fixes its signed path value. Therefore the rank
of the restriction map is exactly the rank $d$ of the observed core rows.
Rank--nullity gives $r_B-d=\rho_{Bj}$, proving
\eqref{eq:observed-variable-count}.
\end{proof}

For clarity, the final block inequalities are
\begin{align}
 y_jL_p&\le\phi_{pj}\le y_jU_p,
       &&j\in J_B,\ p\in\mathcal P_B,\label{eq:observed-state-bounds}\\
 q_{ej}&=\phi_{pj},
       &&(e,j)\in\Obs,\ e\in p\subseteq E_B,\label{eq:observed-equalities}\\
 \lambda_{B,*}L_p&\le s_p(x)-\sum_{j\in J_B}\phi_{pj}
                       \le\lambda_{B,*}U_p,
       &&p\in\mathcal P_B.\label{eq:observed-residual-bounds}
\end{align}
The original flow and simplex domains and bridge observations remain part
of the formulation. In particular, reconstructing state products by balance
equations alone is insufficient: all displayed state and residual bounds
must also be imposed.

The row counts do not bound matrix nonzeros when ranks grow. If
$N=|V|+|E|+m+|\Obs|$ and $\Obs_B=\Obs\cap(E_B\times[m])$, both
constructions have the explicit nonzero bound
\begin{equation}\label{eq:compression-nonzeros}
 O\left(N+\sum_{B:a_B>0}
       \bigl[r_B^2(a_B+1)+r_B|\Obs_B|\bigr]\right).
\end{equation}
Indeed, a reconstructed path has $O(r_B)$ selected-product and auxiliary
terms per label; there are $O(r_B)$ paths and $a_B$ explicit labels.
Each remaining observation also uses $O(r_B)$ terms. Hence bounded maximum
block rank gives size linear in the sparse input, even as $m$ and the graph
grow. The whole formulation need not be totally unimodular.
The unit coefficient statements use the displayed rational row scaling;
clearing denominators in the simplex coefficients can destroy that unit
bound in a primitive integer normalization.

\subsection{Forest complements and product completion}\label{subsec:forest-complements}
\begin{corollary}[An original-coordinate hull on arbitrary graphs]\label{cor:forest-hull}
Suppose $(V_B,U_{Bj})$ is a forest for every $B$ and $j\in J_B$.
Then \eqref{eq:observed-state-bounds}--\eqref{eq:observed-residual-bounds},
the original domains, and the bridge equations describe $\Hull$ with no
additional variables. They have unit flow and product coefficients and
the row bound \eqref{eq:compression-row-count}.
\end{corollary}
\begin{proof}
An undirected multigraph is a forest exactly when its cycle rank is zero.
Thus every $\rho_{Bj}$ in \eqref{eq:observed-variable-count} vanishes.
Apply \cref{thm:observed-compression}.
\end{proof}

Equivalently, the observations for each locally observed label meet every
cycle of its block. Entirely unobserved labels impose no condition: they
are already included in the merged state. The condition is sufficient;
it does not assert that remaining unobserved cycles prevent an
original-coordinate description with unit coefficients.

\begin{proposition}[Minimum individual-coordinate completion]\label{prop:minimum-completion}
Fix $B$ and $j\in J_B$. Among additions of individual missing arc-product
coordinates, exactly $\rho_{Bj}$ are necessary and sufficient to determine
the full state circulation from the observed coordinates and linear balance
equations on the ambient circulation space. A minimum choice consists of
the nonforest edges of any spanning forest of $(V_B,U_{Bj})$.
The additional variables in \cref{thm:observed-compression} may instead be
chosen to represent these individual missing products, with the same
variable and row bounds and unit flow, product, and auxiliary coefficients.
\end{proposition}
\begin{proof}
The existing restriction map has nullity $\rho_{Bj}$. Each additional
scalar coordinate can raise its rank by at most one, proving necessity.
Let $F_{Bj}$ be a spanning forest of the unobserved subgraph, including
all its vertices. Add observations on $U_{Bj}\setminus F_{Bj}$, which
has $\rho_{Bj}$ edges. The remaining unobserved graph is a forest, so
the restriction kernel is zero by the proof of
\cref{thm:observed-compression}. This proves sufficiency and the stated
minimum.

Apply \cref{cor:forest-hull} to the completed observation set and regard
the added product coordinates as auxiliary variables. Projection onto the
original observations commutes with convexification, so this is an extended
formulation for the original $\Hull$. Completing labels only in $J_B$
does not introduce new labels. The added equations number
$\sum_{B,j}\rho_{Bj}\le\sum_Br_Ba_B$, already absorbed by the row
bound. The corollary's coefficient property includes the added coordinates.
\end{proof}

There is a direct interpretation of this completion. The block state flows
satisfy $A_Bf^{B,j}=y_jA_Bv_{E_B}$, since their deviations are block
circulations. Once the nonforest state products are known, sum these
equations on either side of a cut formed by deleting a remaining forest
edge from its component. That edge is the only
unknown crossing edge, so its state product is an affine expression in the
known state products and the state weight, with unit product coefficients.
Component balance equations enforce consistency if the unobserved forest
is disconnected. This is network elimination of the kind underlying reduced
RLT: \citet[Theorem~3.1 and Section~2]{LibertiPantelides2006} reconstruct
common-factor products by multiplying full-row-rank linear equations and
retaining products on suitable complementary coordinates. Here the extra
information is the observation-complement graph criterion and its exact
sparse simplex-hull consequence.

The minimum in \cref{prop:minimum-completion} concerns reconstruction by
\emph{individual coordinates on the ambient linear space}. It is neither
an extension-complexity lower bound nor a claim about the number of
coordinates needed on a capacity-degenerate slice. In particular,
$\lambda_j=0$ forces every state product to zero. The preprocessing in
\cref{lem:fixed-arcs} can remove fixed arcs before the parameters are
computed, but it does not remove this conditional-slice distinction.

All formulations also give explicit recovery. From retained coordinates,
recover each $t^{B,j}$ using \eqref{eq:pivot-reconstruction}, reconstruct
path deviations by $\widehat C_Bt^{B,j}$, and subtract their sum from
$s(x)$ to recover the merged state. Expand deviations along the original
paths using their signs. The proportional refinement
\eqref{eq:proportional-refinement} then supplies the global state flows
and the at-most-$m+1$-point decomposition from
\cref{prop:disaggregation}. This recovery is exact rational arithmetic for
rational input and a rational feasible extension; merely solving the
formulation numerically does not provide an exact certificate.

\begin{example}[Three observations on a $K_4$ block]\label{ex:forest-k4}
Orient every edge of $K_4$ from the smaller to the larger vertex and take
unit capacities and unit flow from vertex $1$ to vertex $4$. Thus
$b=(-1,0,0,1)^\top$. Use the reference flow $v_{14}=1$ and zero on the
other arcs. For one label $j$, observe $13$, $14$, and $24$; the unobserved
edges $12$, $23$, and $34$ form a spanning tree. Writing
$q_{13}=z_{13,j}$, $q_{14}=z_{14,j}-y_j$, and $q_{24}=z_{24,j}$, balance
gives the remaining state deviations
\[
 h_{12}=-q_{13}-q_{14},\qquad
 h_{23}=-q_{13}-q_{14}-q_{24},\qquad
 h_{34}=-q_{14}-q_{24}.
\]
For example, nonnegativity of the reconstructed flow on $23$ gives
\begin{equation}\label{eq:k4-forest-cut}
 z_{13,j}+z_{14,j}+z_{24,j}\le y_j.
\end{equation}
This cut is absent from independent product convexification. To see the
gap, take $m=1$, $y_1=1/2$, $z_{13,1}=z_{14,1}=z_{24,1}=1/5$, and let
$x$ average the four paths $1234$, $124$, $134$, and $14$ equally. Then
$x_{12}=x_{34}=1/2$ and the other four arc flows are $1/4$. All three
observed products satisfy \eqref{eq:mccormick}, but the left side of
\eqref{eq:k4-forest-cut} is $3/5>1/2$. The complete hull requires all
reconstructed bounds, including the residual bounds, not only this one cut.
\end{example}


\bibliographystyle{plainnat}
\bibliography{references}
\end{document}
